\documentclass[12pt,a4paper]{article}

\usepackage{fontawesome5}
\providecommand{\faPencilAlt}{\faIcon{pencil-alt}}

\usepackage{styles/couleurs}
\usepackage{styles/mise_en_page}
\usepackage{styles/boites}
\usepackage{styles/maths}
\usepackage{styles/figures}

\usepackage{enumitem}
\usepackage{array}
\usepackage{tabularx}
\usepackage{amsmath}
\usepackage{amssymb}
\usepackage{graphicx}
\usepackage{tikz}
\usepackage{xparse}

\usetikzlibrary{calc,angles,quotes,arrows.meta}

\setlist[itemize]{leftmargin=1.4em,itemsep=0.15em,topsep=0.2em}
\setlist[enumerate]{leftmargin=*,itemsep=0.22em,topsep=0.25em}

\newcommand{\lignepointillee}{\noindent\dotfill\par\vspace{0.45em}}

\newcommand{\lignesreponse}[1]{%
\begin{tcolorbox}[
enhanced,colback=white,colframe=NeutreC!40,boxrule=0.5pt,
arc=1.5mm,left=3mm,right=3mm,top=2mm,bottom=1mm,
before skip=2pt,after skip=2pt]
\foreach \i in {1,...,#1}{\lignepointillee}
\end{tcolorbox}}

\newcommand{\titreseq}[1]{%
\begin{center}
{\sffamily\large\bfseries\color{SavoirC}#1}
\end{center}}

\begin{document}

\titreseq{Activité 3 -- Nommer avant de calculer}

\section*{Partie A -- Le côté adjacent dépend de l'angle}

Dans le triangle \(ABC\) rectangle en \(B\) :

\begin{center}
\begin{tikzpicture}[scale=1]
  \coordinate (A) at (0,0);
  \coordinate (B) at (5,0);
  \coordinate (C) at (5,3.5);
  \draw[line width=0.9pt] (A)--(B)--(C)--cycle;
  \draw (4.65,0)--(4.65,0.35)--(5,0.35);
  \draw pic[draw=ActionC,fill=ActionC!15,angle radius=8mm] {angle=B--A--C};
  \draw pic[draw=SavoirC,fill=SavoirC!12,angle radius=7mm] {angle=A--C--B};
  \node[below left] at (A) {$A$};
  \node[below right] at (B) {$B$};
  \node[above right] at (C) {$C$};
\end{tikzpicture}
\end{center}

Complète.

\renewcommand{\arraystretch}{1.55}
\begin{tabularx}{0.88\textwidth}{
|>{\raggedright\arraybackslash}p{4.2cm}
|>{\centering\arraybackslash}X
|>{\centering\arraybackslash}X|}
\hline
 & \textbf{Angle en \(A\)} & \textbf{Angle en \(C\)}\\
\hline
\textbf{Hypoténuse} & & \\
\hline
\textbf{Côté adjacent} & & \\
\hline
\end{tabularx}

\vspace{0.3cm}

Explique pourquoi l'hypoténuse ne change pas, mais le côté adjacent change.

\vspace{0.3cm}


\lignesreponse{5}

\section*{Partie B -- Définition du cosinus}

Dans un triangle rectangle, pour un angle aigu donné :

\[
\cos(\text{angle})
=
\text{côté adjacent}\div\text{hypoténuse}.
\]

\vspace{0.3cm}

Complète dans le triangle précédent :

\[
\cos(\widehat{A})
=
\ldots\div\ldots
\]

\[
\cos(\widehat{C})
=
\ldots\div\ldots
\]

\newpage

\section*{Partie C -- Un protocole qui évite les confusions}

Remets les étapes dans l'ordre.

\begin{itemize}
  \item[\(\square\)] Je repère l'angle aigu étudié.
  \item[\(\square\)] Je nomme le triangle rectangle dans lequel je travaille.
  \item[\(\square\)] Je repère le côté adjacent à cet angle.
  \item[\(\square\)] Je repère l'hypoténuse, opposée à l'angle droit.
  \item[\(\square\)] J'écris le cosinus avec les noms des segments.
\end{itemize}

\section*{Partie D -- La calculatrice en mode degré}

Vérifie que la calculatrice est en mode degré, puis complète au centième.

\[
\cos 20^\circ\approx\ldots
\qquad
\cos 35^\circ\approx\ldots
\]

\[
\cos 60^\circ=\ldots
\qquad
\cos 80^\circ\approx\ldots
\]

Que remarques-tu lorsque l'angle aigu augmente ?

\lignesreponse{4}

\aretenir{Le cosinus d'un angle aigu est un nombre sans unité compris entre \(0\) et \(1\).}

%==================================================
% Encart calculatrice : mode degré
%==================================================

\begin{tcolorbox}[
  enhanced,
  colback=SavoirBg,
  colframe=SavoirC!65,
  boxrule=0.7pt,
  arc=2mm,
  left=3mm,
  right=3mm,
  top=2mm,
  bottom=2mm,
  title={\sffamily\bfseries
    \faCalculator\ Calculatrice : vérifier le mode degré},
  coltitle=SavoirC,
  fonttitle=\sffamily\bfseries
]

Avant d'effectuer un calcul avec un cosinus, vérifie que la calculatrice
est réglée en \textbf{degrés}.

\vspace{0.3em}

\renewcommand{\arraystretch}{1.35}
\begin{tabularx}{\linewidth}{
  @{}
  >{\sffamily\bfseries\color{TitreC}}p{3.2cm}
  X
  @{}
}
\textbf{Casio fx-92}
&
\texttt{CONFIG}
\(\longrightarrow\)
\texttt{Unité d'angle}
\(\longrightarrow\)
\texttt{Degré}
\\

\textbf{TI-Collège Plus}
&
\texttt{MODE}
\(\longrightarrow\)
sélectionner \texttt{DEG}
\(\longrightarrow\)
valider
\end{tabularx}

\begin{center}
  \sffamily\bfseries
  L'écran doit afficher
  \tcbox[
    on line,
    colback=white,
    colframe=SavoirC,
    boxrule=0.6pt,
    arc=1mm,
    left=1.5mm,
    right=1.5mm,
    top=0.4mm,
    bottom=0.4mm
  ]{\texttt{D} ou \texttt{DEG}}.
\end{center}

\textbf{Test rapide :}

\[
\cos(60^\circ)=0{,}5.
\]

Si la calculatrice n'affiche pas \(0{,}5\), vérifie à nouveau le mode.

\end{tcolorbox}

\end{document}
